The curvature statements below are computed in Gauge Theory Gravity on flat spacetime (Definition \claimref{obj.synthesis.gtg-flat-spacetime-gauge-fields}): each metric is, up to the overall sign of dictionary C7 (the line elements are written in $(-,+,+,+)$), the derived metric of a diagonal gauge field $\hb^{-1}(e_{t}) = N(t)\,e_{0}$, $\hb^{-1}(e_{x^{i}}) = A_{i}(t)\,e_{i}$ with constant orthonormal physical frame $e_{0}, \dots, e_{3}$ and $B_{i} = e_{0}\wedge e_{i}$. For such a field the torsion equation gives $\Omega(e_{t}) = 0$, $\Omega(e_{x^{i}}) = -(\dot A_{i}/N)\,B_{i}$, so $R(e_{t}\wedge e_{x^{i}}) = -\partial_{t}(\dot A_{i}/N)\,B_{i}$, $R(e_{x^{i}}\wedge e_{x^{j}}) = -(\dot A_{i}\dot A_{j}/N^{2})\,e_{i}\wedge e_{j}$, and the store-sign operator is diagonal on the frame planes:
\[
\mathcal{R}(B_{i}) = \frac{\partial_{t}(\dot A_{i}/N)}{N A_{i}}\,B_{i}, \qquad \mathcal{R}(e_{i}\wedge e_{j}) = \frac{\dot A_{i}\dot A_{j}}{N^{2}A_{i}A_{j}}\,e_{i}\wedge e_{j}.
\]
On a plane with eigenvalue $\kappa$ the contraction $\partial_{a}\cdot\mathcal{R}(a\wedge b)$ contributes $\kappa$ to each of its two vectors, so the Ricci map is diagonal with $\mathcal{R}(e_{0}) = \sum_{i}\kappa_{0i}\,e_{0}$, $\mathcal{R}(e_{j}) = \bigl(\kappa_{0j} + \sum_{i \neq j}\kappa_{ij}\bigr)e_{j}$, and the scalar is $\mathcal{R} = 2\sum_{\text{planes}}\kappa$. The values below are store-sign values in $(+,-,-,-)$; the scalar and the Ricci square have the same values in the check's index convention in $(-,+,+,+)$, as the statement records ($g^{\circ} \to -g^{\circ}$ leaves the Ricci tensor of the index convention unchanged and flips the inverse metric, and the store's Ricci tensor is its negative, so the scalar is the same). The check computes each field with the library's engine on exact jets, $\Omega$ solved from the torsion equation (check\_gtg\_wall\_field\_strengths).

\emph{The models are tame.} For $G_{a,b} = \mathrm{diag}(-t^{a}, t^{b}, 1, 1)$, $\det G_{a,b} = -t^{a + b}$, so the degeneration at $t = 0$ is of corank one and first order exactly when $a + b = 1$; the pairs $(2, 0)$, $(0, 2)$, $(1, 1)$ are not tame. For $(1, 0)$ the radical is spanned by $k = dt$, the conormal of $\Sigma$, with $G(k, k) = -t$ and transversality datum $\partial_{t}G(k, k) = -1 \neq 0$; for $(0, 1)$ it is spanned by $k = dx$, tangent to $\Sigma$ in the sense that $k$ is not a multiple of $dt$, with $G(k, k) = t$. Both satisfy (T1) to (T3) of Definition \claimref{sheaf.tameness}, and in both the metric has a simple pole in one direction ($g_{tt} = -t^{-1}$, respectively $g_{xx} = t^{-1}$), as Fact \claimref{dark-energy.imported-walls-sphere-theorem} and Theorem \claimref{sheaf.metrics-residue-square-root-twist} require.

\emph{(i)} With $\tau = 2\sqrt{t}$, $dt = \tfrac{\tau}{2}\,d\tau$ and $t = \tfrac{\tau^{2}}{4}$, so $-t^{-1}dt^{2} = -d\tau^{2}$ and $g = -d\tau^{2} + dx^{2} + dy^{2} + dz^{2}$ on $t > 0$: the half $\tau > 0$ of Minkowski space. The proper time from $\Sigma$ to the slice $t$ is $\int_{0}^{t}s^{-1/2}ds = 2\sqrt{t}$. On the half space every inextendible causal curve has $\tau$ strictly monotone, tending to $0$ in the past and to $\infty$ in the future, so it meets each slice $\tau = c > 0$ exactly once: the slices are Cauchy surfaces of the physical side. For $t < 0$ the cometric is $\mathrm{diag}(|t|, 1, 1, 1)$, positive definite. In Gauge Theory Gravity the field is $N = t^{-1/2}$, $A_{i} = 1$: each $\hb^{-1}(e_{\mu})$ depends only on its own label, so $\Omega = 0$ solves the torsion equation and the field strength vanishes: the physical side is flat in the coordinate $t$ as well (the check's coordinate engine also returns vanishing Ricci tensor). Since each $\tau = \varepsilon$ is a Cauchy surface, the propagation of the Hadamard form (Fact \claimref{synthesis.imported-propagation-of-hadamard-form}, scalar field) carries a two-point function of Hadamard form on a neighbourhood of $\tau = \varepsilon$ to the whole physical side.

\emph{(ii), scalar.} $\partial_{\tau} = \tfrac{dt}{d\tau}\partial_{t} = \sqrt{t}\,\partial_{t}$, hence $\partial_{\tau}^{2} = \sqrt{t}\,\partial_{t}(\sqrt{t}\,\partial_{t}) = t\,\partial_{t}^{2} + \tfrac12\partial_{t}$, whose indicial polynomial on $t^{r}$ is $r(r - 1) + \tfrac{r}{2} = r(r - \tfrac12)$, with exponents $0$ and $\tfrac12$. A solution of $(-\partial_{\tau}^{2} + \Delta - m^{2})\phi = 0$ on the physical side is the restriction of a solution on Minkowski space, smooth in $\tau$ at $\tau = 0$, and the reflection $\tau \mapsto -\tau$ maps solutions to solutions, so $\phi = \phi_{+} + \phi_{-}$ with $\phi_{\pm}$ even and odd in $\tau$. A smooth even function of $\tau$ is a smooth function of $\tau^{2} = 4t$ (Whitney's theorem on even functions, Fact \claimref{synthesis.imported-bochner-theorem-moment-lemma-and-whitney-even-functions}(c)), so $\phi_{+}$ is smooth in $t$ up to $\Sigma$. If $\phi_{-} \neq 0$, its Cauchy data at $\tau = 0$ are $(0, \partial_{\tau}\phi_{-})$ with $\partial_{\tau}\phi_{-} \neq 0$ (zero data give the zero solution), so $\phi_{-} = \tau\,\psi(\tau^{2})$ with $\psi(0) \neq 0$, that is $2\sqrt{t}\,\psi(4t)$, whose $t$-derivative is unbounded at $t = 0$; and $\phi_{+}$ is smooth in $t$, so $\phi$ is smooth in $t$ up to $\Sigma$ exactly when $\phi_{-} = 0$, that is when $\partial_{\tau}\phi = 0$ on $\Sigma$. In the coordinate $t$ the conjugate momentum is $\sqrt{t}\,\partial_{t}\phi$, which vanishes on $\Sigma$ for every $\phi$ smooth in $t$. On solutions with data $(f, 0)$ the symplectic form $\sigma(\phi_{1}, \phi_{2}) = \int_{\Sigma}(\phi_{1}\partial_{\tau}\phi_{2} - \phi_{2}\partial_{\tau}\phi_{1})\,d^{3}x$ vanishes, and they are maximal with this property, since $\sigma((f, 0), (0, g)) = \int fg$ pairs them nondegenerately with the data $(0, g)$: a Lagrangian subspace. The field operators smeared with these solutions therefore commute, and no field obeying the canonical commutation relations on the physical side is smooth in $t$ across $\Sigma$. The check verifies the mode form exactly: for $\omega^{2} = k^{2} + m^{2}$ at three rational values, $\cos(\omega\tau)$ written as a series in $t$ satisfies $(t\,\partial_{t}^{2} + \tfrac12\partial_{t})C = -\omega^{2}C$ coefficient by coefficient (to order $t^{6}$), $\sin(\omega\tau)/\omega$ has only odd powers of $\tau$ (the exponent $\tfrac12$ branch), $\sigma(\cos, \cos) = 0$ and $\sigma(\cos, \sin/\omega) = 1$ at $\tau = 0$, and $\sqrt{t}\,\partial_{t}C$ starts at $t^{1/2}$.

\emph{(ii), Dirac.} Work in the constant frame of the chart $(\tau, x, y, z)$, in which the spin connection vanishes. In Hestenes form $\nabla\psi I\sigma_{3} = m\psi\gamma_{0}$ with $\nabla = \gamma^{0}\partial_{\tau} + \gamma^{i}\partial_{i}$. Put $\psi = \phi(\tau)\,e$, $e = \exp(-I\sigma_{3}\,k\cdot x)$, $\phi$ in the even subalgebra. Since $I\sigma_{3} = -\gamma_{1}\gamma_{2}$ commutes with $\gamma_{0}$ and $(I\sigma_{3})^{2} = -1$, one has $\partial_{i}e = -k_{i}\,e\,I\sigma_{3}$, and the equation becomes, after cancelling the invertible factor $e$ on the right, $\gamma_{0}\phi' I\sigma_{3} + \kappa\phi = m\phi\gamma_{0}$ with $\kappa = k_{i}\gamma^{i}$. So $\phi' = M\phi$, $M\phi = \gamma_{0}(m\phi\gamma_{0} - \kappa\phi)(I\sigma_{3})^{-1}$, a linear map of the even subalgebra (the check uses $\kappa = \sum k_{i}\gamma_{i}$; the sign of $k$ is immaterial). The check computes, exactly on the eight even blades and five samples of $(k, m)$, that $M$ preserves the even subalgebra and $M^{2} = -\omega^{2}$, $\omega^{2} = |k|^{2} + m^{2}$; by hand, write $J = I\sigma_{3}$ and $L_{a}$, $R_{b}$ for left and right multiplication, which commute, with $R_{b}R_{c} = R_{cb}$; then $M = A + B$, $A = m L_{\gamma_{0}}R_{\gamma_{0}J^{-1}}$, $B = -L_{\gamma_{0}\kappa}R_{J^{-1}}$, and since $J$ commutes with $\gamma_{0}$, $J^{2} = -1$, $\gamma_{0}\kappa = -\kappa\gamma_{0}$ and $\kappa^{2} = -|k|^{2}$: $A^{2} = m^{2}R_{(\gamma_{0}J^{-1})^{2}} = -m^{2}$, $B^{2} = L_{(\gamma_{0}\kappa)^{2}}R_{J^{-2}} = -|k|^{2}$, and $AB + BA = -m(L_{\kappa} - L_{\kappa})R_{-\gamma_{0}} = 0$, so $M^{2} = -(m^{2} + |k|^{2})$. Hence $\phi(\tau) = \cos(\omega\tau)\phi_{0} + \omega^{-1}\sin(\omega\tau)M\phi_{0}$, whose odd part vanishes exactly when $M\phi_{0} = 0$, and $M$ is invertible for $\omega \neq 0$ (the check finds rank $8$, and rank $0$ at $k = 0 = m$). By the argument of the scalar case a Dirac plane wave is smooth in $t$ across $\Sigma$ only if it is even in $\tau$, so only if $\phi_{0} = 0$, unless $k = 0 = m$. A solution with square-integrable data on the flat slice is a superposition of plane waves whose odd part is $\int \omega_{k}^{-1}\sin(\omega_{k}\tau)M_{k}\phi_{0}(k)\,d^{3}k$; it vanishes exactly when $M_{k}\phi_{0}(k) = 0$ for almost every $k$, that is when $\phi_{0} = 0$. So the only Dirac solution with square-integrable data that is smooth in $t$ across $\Sigma$ is zero, for every mass. The frame matters: a spin rotation odd in $\tau$, for example $R = \exp(-I\sigma_{3}m\tau)$ at $k = 0$, makes the rest-frame solution constant, at the price of a spin connection $\Omega(\partial_{t}) \propto t^{-1/2}$, not smooth in $t$; the statement is for the constant frame, or equivalently for gauge fields $(\hb, \Omega)$ whose spin connection vanishes there. \emph{(Argued, beyond the model.)} With a transverse metric $h(t)$ smooth in $t$, the spin-connection term along $\gamma^{0}$ is odd in $\tau$ and the transverse Dirac operator with mass, $K(\tau)$, is even, so an even solution must satisfy $K\psi = 0$ and is a static zero mode of the transverse Dirac Hamiltonian, of which there is none in $L^{2}$ when $m \neq 0$, since that Hamiltonian squares to $D_{h}^{2} + m^{2} \geq m^{2}$.

\emph{(iii)} For $(0, 1)$, $g = -dt^{2} + t^{-1}dx^{2} + dy^{2} + dz^{2}$ and the proper time to $\Sigma$ is $t$. The field is $N = t^{-a/2}$, $A_{1} = t^{-b/2}$, $A_{2} = A_{3} = 1$: only the plane $B_{1}$ carries field strength, with $\kappa = \partial_{t}(\dot A_{1}/N)/(NA_{1})$; $\dot A_{1}/N = -\tfrac{b}{2}t^{(a - b)/2 - 1}$, so $\kappa = \tfrac{b}{4}(b + 2 - a)\,t^{a - 2}$ and $\mathcal{R} = 2\kappa = \tfrac{b}{2}(b + 2 - a)\,t^{a - 2}$; for $(0, 1)$ ($N = 1$, $\kappa = \ddot A_{1}/A_{1}$, $A_{1} = t^{-1/2}$, $\ddot A_{1} = \tfrac34 t^{-5/2}$) this is $\mathcal{R} = \tfrac32 t^{-2}$. The check computes it with the engine for all five pairs, and its coordinate engine ($R^{\rho}{}_{\sigma\mu\nu} = \partial_{\mu}\Gamma^{\rho}_{\nu\sigma} - \partial_{\nu}\Gamma^{\rho}_{\mu\sigma} + \Gamma^{\rho}_{\mu\lambda}\Gamma^{\lambda}_{\nu\sigma} - \Gamma^{\rho}_{\nu\lambda}\Gamma^{\lambda}_{\mu\sigma}$, $R = 2\ddot A/A$ for $-dt^{2} + A(t)^{2}dx^{2}$ plus a flat block) agrees. For $t < 0$ the cometric $\mathrm{diag}(-1, -|t|, 1, 1)$ has two negative entries: signature $(2, 2)$. On $t > 0$, $t$ is a time function and causal curves have $|dx/dt| \le \sqrt{t}$, so each inextendible one meets every slice $t = \varepsilon$ once: these slices are Cauchy surfaces, and propagation acts on $t > 0$ from any of them; at $\Sigma$ itself the curvature is unbounded, and a state there is initial data to be matched, as in Premise \claimref{synthesis.pre-descent-state-carried-into-descended-phase-premise}.

\emph{(iv)} $-dt^{2} + t^{2p}\delta = t^{2p}(-t^{-2p}dt^{2} + \delta)$; the bracket's cometric $\mathrm{diag}(-t^{2p}, 1, 1, 1)$ has radical $dt$ at $t = 0$ for $p > 0$ and determinant of order $2p$, so it is tame exactly when $2p = 1$. For a spatially flat field of a perfect fluid with $P = w\rho$ the standard solution is $a = t^{p}$, $p = 2/(3(1 + w))$; in Gauge Theory Gravity the field $N = 1$, $A_{i} = t^{p}$ has $\kappa_{0i} = p(p - 1)t^{-2}$, $\kappa_{ij} = p^{2}t^{-2}$, so $\mathcal{R}(e_{0}) = 3p(p - 1)t^{-2}e_{0}$, $\mathcal{R}(e_{j}) = p(3p - 1)t^{-2}e_{j}$, $\mathcal{R} = 6p(2p - 1)t^{-2}$, and the Einstein map is $\mathcal{G}(e_{0}) = -3p^{2}t^{-2}e_{0}$, $\mathcal{G}(e_{j}) = p(2 - 3p)t^{-2}e_{j}$; the field equation $\mathcal{G}(a) = -8\pi T(a)$ with $T(e_{0}) = \rho e_{0}$, $T(e_{j}) = -Pe_{j}$ gives $\rho = 3p^{2}/(8\pi t^{2})$, $P = p(2 - 3p)/(8\pi t^{2})$, and $P = w\rho$ is $2 - 3p = 3wp$, $p = 2/(3(1 + w))$. The check recomputes it at five values of $w$ (in GTG and with the coordinate engine, $G_{xx} = w\,a^{2}G_{tt}$) and finds the tame case only at $w = \tfrac13$, $p = \tfrac12$. Then $g_{\mathrm{FRW}} = t\,g_{1,0}$, and the wall's $\tau = 2\sqrt{t}$ is the conformal time $\int dt/a$. At $p = \tfrac12$: $\mathcal{R}(e_{0}) = -\tfrac34 t^{-2}e_{0}$, $\mathcal{R}(e_{j}) = \tfrac14 t^{-2}e_{j}$, so $\mathcal{R} = 0$ and the Ricci square is $\mathcal{R}(\partial_{a})\cdot\mathcal{R}(a) = \sum_{a}\mathcal{R}(e^{a})\cdot\mathcal{R}(e_{a}) = \tfrac{9}{16}t^{-4} + 3\cdot\tfrac{1}{16}t^{-4} = \tfrac34 t^{-4}$ (in components $R_{ab}R^{ab}$, which the coordinate engine also returns). The massless Dirac field and the Maxwell and Yang-Mills fields are conformally covariant in four dimensions (Fact \claimref{synthesis.imported-conformal-covariance-of-massless-fields-and-local-conformal-flatness}(b); the stored Fact \claimref{dark-energy.gravity-pre-metric-forces-degenerate-exactly} names the Yang-Mills actions as conformal-level objects), so for them the radiation-dominated field is, up to the factor $t$, the physical side of (i): the origin is the initial slice $\tau = 0$ of half of Minkowski space, a Cauchy surface of its extension, and the Weyl curvature vanishes there. The conformal-time reading is not special to radiation: for every $p < 1$ the conformal time $\eta = t^{1 - p}/(1 - p)$ is finite, $d\eta/dt = t^{-p}$ turns the bracket into $-d\eta^{2} + \delta$ (the check verifies $-t^{-2p}(dt/d\eta)^{2} = -1$ at the four sampled $w$ with $p < 1$), and in $t' = \eta^{2}/4$ this is $g_{1,0}$, by the computation of (i). What singles out $w = \tfrac13$ is only that the bracket is a tame wall in the cosmic coordinate $t$. \emph{Smoothness in $t$ at the radiation origin depends on the conformal weight, not on (ii).} With $a = t^{1/2} = \tau/2$, a conformally covariant field of weight $s$ is $a^{-s}$ times a solution on the flat half-space. For the massless scalar ($R = 0$, so minimal and conformal coupling agree) $\phi = t^{-1/2}\phi_{\mathrm{flat}}$; its equation $\phi'' + \tfrac{3}{2t}\phi' + \tfrac{k^{2}}{t}\phi = 0$ has indicial exponents $0$ and $-\tfrac12$, and the solutions smooth in $t$ are those with $\phi_{\mathrm{flat}}$ odd in $\tau$ (for example $2\sin(k\tau)/\tau$, which the check expands in $t$ and verifies coefficient by coefficient at three values of $k^{2}$, while $2\cos(k\tau)/\tau$ starts at $t^{-1/2}$): the opposite half from (ii). For the massless Dirac field $\psi = a^{-3/2}\psi_{\mathrm{flat}} = t^{-3/4}\psi_{\mathrm{flat}}$, and $\psi_{\mathrm{flat}}$ smooth in $\tau$ contributes powers $\tau^{n} = 2^{n}t^{n/2}$, so the exponents of $\psi$ are $-\tfrac34 + \tfrac{n}{2}$, never integers: no nonzero solution is smooth in $t$, whatever its parity. So which solutions are smooth in $t$ is weight-dependent, and at the radiation origin the cosmic coordinate supplies no Dirac data for a reason different from (ii). \emph{(Argued.)} For a gauge potential smooth in $t$ the electric field $F_{\tau i} = \sqrt{t}\,F_{ti}$ vanishes on $\Sigma$: the same Lagrangian half.

\emph{Reading (argued, not part of the statement).} The chain of op.deposit-state-is-hadamard needs, at its second step, the descent's surface to be a Cauchy surface on which the deposit is Hadamard. At an aligned tame wall with the radical normal, and for the conformal fields at a radiation-dominated origin, the surface is a regular initial slice in the proper (conformal) time and the propagation step holds from any collar $0 < \tau < \varepsilon$; but smoothness at the wall supplies no state: in its own smooth structure it carries a Lagrangian half of the bulk scalar data and none of the Dirac data in the constant frame, and in the proper time it is a regular boundary imposing no condition. Two possible selectors are not examined here: the boundary modes that the wall carries by Fact \claimref{dark-energy.imported-walls-sphere-theorem} (the Witt partner and its bilinears), which are not bulk fields, and the Riemannian far side read as the domain of a Euclidean (no-boundary) path integral, which for the flat half-space would give the vacuum's wave functional. The descent's conditional expectation selects nothing either: it is the identity on its range (Theorem \claimref{transport.pinching-kernel}(i)), so there it returns the state it is applied to, and outside its range it only erases the coherences with the adjoined pair; every state, Hadamard or not, is a possible input. The missing datum is therefore a Hadamard state on a collar of the wall, which neither smoothness at the wall nor the descent's expectation selects. At a wall with the radical tangent, and at the Kasner walls of (P1) and (P2), the surface is singular and the state there is initial data to be matched. Which case the descent is, the store does not fix: the dating's pre-metric reading (\claimref{standard-model.dating-descent}, a conjecture) is silent on the wall's normal form, and the far side of a radical-normal wall is Riemannian rather than pre-metric. The propagation reaches only the domain of dependence of the collar, by the definition of a Cauchy surface: where a later field has a Cauchy horizon, as the store's charged and rotating holes do, it stops there; for the Reissner-Nordstr\"om-de Sitter field the renormalised stress tensor of a Hadamard state generically diverges at the Cauchy horizon (Hollands, Wald and Zahn 2020, Fact \claimref{synthesis.imported-hadamard-stress-tensor-diverges-at-the-reissner-nordstrom-de-sitter-cauchy-horizon}). Transport-absent: the computations are standard signature-change and conformal geometry, applied to the series' tame walls. Signature change at a degenerate metric, and the question whether momenta vanish at the change, were studied by Ellis, Sumeruk, Coule and Hellaby (Class. Quantum Grav. 9 (1992) 1535) and by Dray, Manogue and Tucker (Gen. Rel. Grav. 23 (1991) 967), cited from memory; there the metric, not the cometric, degenerates.
